\addchap{Abstract}

This project addresses the next mandatory steps bridging from a first project on the topic towards the master's thesis. The problem addressed is the identification of a suitable tool with which the automation of the current data quality checks can be performed. During the course of this work, another main goal is to provide the business with tools that can assist them with the usage of new data sources. The Power BI report for the daily checks of the message processing data did exactly that, employing improved visualisations developed by the project team. All this was done while applying the DevOps principles at Swiss Post and implementing user feedback received while building the report for stakeholders.

As a next step, the search for a DQM tool was addressed by following the team during the setup of requirement lists and the evaluation of tools, including demonstrations of such tools. However, as the decision on the tool had not yet been taken by the end of this project, the way forward was defined together with the DWH LS OPS solution architect in a meeting. The monitoring that was set up to check the uptime of jobs such as ETL jobs was selected for testing in its second application, namely running checks on data. This tool was able to run counts on queries executed through Amazon Athena. However, the function had not yet been properly tested. Therefore, the project team carried out tests to verify whether this tool could be applied to the use case defined for the master's thesis. This use case was the automation of the data quality checks performed daily, which were part of the previous research.

With this information, the next step was to test this possible data quality tool, which was developed during the AWS Modernisation and was part of the full BizDevOps team until the end of May. Until then, knowledge transfer sessions were performed. After these transfers, a test was derived based on the full use case of the master's thesis, reducing its complexity and amount of data to perform only one count of one specific sorting centre for one event rather than for every sorting centre.

Finally, it was possible to evaluate the given solution as suitable for the master's thesis. This included not only the setup of a test case but also the setup of the required infrastructure. Thus, everything required for the master's thesis has been tested successfully. The details of this research, the additions to the solution landscape, and the implementation are part of the documentation that follows.